\documentclass[a4paper,12pt]{article}
\input{LaTeX-ToolKit/template.tex}
\begin{document}
\title{Japanese}
\author{沈威宇}
\date{\temtoday}
\titletocdoc
\begin{longtable}[c]{|C{0.2\textwidth}|C{0.1\textwidth}|C{0.1\textwidth}|C{0.1\textwidth}|C{0.1\textwidth}|C{0.1\textwidth}|}
\hline
 & 一段（いちだん）活用（かつよ）動詞（どうし） & イ形容詞（けいようし）/形容詞 & イ形容詞良い（いい）(First い or よ corresponds to 漢字良) & ナ形容詞/形容動詞（けいようどうし） & 名詞（めいし） \\\hline
普通形（ふつうけ）非過去形（ひかこけ）肯定（こうて） & 〜る & 〜い & いい & 〜だ & 〜だ \\\hline
丁寧形（ていねいけ）非過去形肯定 & 〜ます & 〜いです & いいです & 〜です & 〜です \\\hline
普通形非過去形否定（ひて） & 〜ない & 〜くない & よくない & 〜じゃない/〜ではない & 〜じゃない/〜ではない \\\hline
丁寧形非過去否定 & 〜ません & 〜くないです/〜くありません & よくないです/よくありません & 〜じゃありません/〜ではありません & 〜じゃありません/〜ではありません \\\hline
普通形過去形（かこけ）肯定 & 〜た & 〜かった & よかった & 〜だった & 〜だった \\\hline
丁寧形過去形肯定 & 〜ました & 〜かったです & よかったです & 〜でした & 〜でした \\\hline
普通形過去形否定 & 〜なかった & 〜くなかった & よくなかった & 〜じゃなかった/〜ではなかった & 〜じゃなかった/〜ではなかった \\\hline
丁寧形過去形否定 & 〜ませんでした & 〜くなかったです/〜くありませんでした & よくなかったです/よくありませんでした & 〜じゃなかったです/〜じゃありませんでした/〜ではなかったです/〜ではありませんでした & 〜じゃなかったです/〜じゃありませんでした/〜ではなかったです/〜ではありませんでした \\\hline
\end{longtable}
五段活用/グループ1
（上/下）一段活用/グループ2
来る
する
\end{document}
